\section{Train Time 与 Test Time：哪些组件真正被训练}

理解 RAG 论文最有效的方法，不是按年份背模型名，而是问两个问题：训练时更新什么，测试时运行什么。Frozen RAG、RePlug、RAG、REALM 和 Atlas 的核心差异，都能映射到 language model、query encoder、document encoder、reranker、index 与 prompt 的更新边界。

\subsection{训练与推理要分开列账}

Train Time and Test Time slide 左侧询问是否更新 LM、query encoder、document encoder、是否一次更新全部，以及是否从零预训练；右侧询问系统是否是 frozen model、retrieval 如何变化、不同例子是否使用不同或同一 index。

\lecturefigure{slide-10-train-test-time.jpg}{RAG 在训练期与推理期的选择矩阵}{Stanford 课堂视频 00:10:54}

\begin{knowledgebox}{读图：同一个组件可能在训练时冻结、推理时仍参与}
``Frozen'' 指参数不更新，不等于组件不运行。Frozen retriever 仍会搜索，frozen generator 仍会生成；index 也可能在不训练 encoder 的情况下更新文档。论文比较时必须同时标注 parameter update、index refresh 与 inference calls。
\end{knowledgebox}

\begin{importantbox}{参数与状态的分层}
令参数集合为
\[
\Theta=\{\theta_G,\theta_Q,\theta_D,\theta_R\},
\]
分别代表 generator、query encoder、document encoder 与 reranker；令 $I_t$ 表示时刻 $t$ 的 index 状态。训练可以更新 $\Theta$ 的子集，部署可以只替换 $I_t$。参数学习与知识库更新是两条不同生命周期。
\end{importantbox}

\subsection{Frozen RAG：不训练也能得到强基线}

Frozen RAG slide 把 no training 与 in-context learning 放在一起：用现成 retriever 或 vector database 找文档，把文档拼进 prompt，再调用冻结 LLM。它实现快、适合 black-box API，但 retriever 和 generator 并未为彼此优化。

\lecturefigure{slide-11-frozen-rag.jpg}{Frozen RAG：retriever 与 LLM 通过 prompt 松耦合}{Stanford 课堂视频 00:11:57}

\begin{knowledgebox}{读图：Frozen RAG 的工程价值与研究局限}
价值在于模块可替换、无需模型权重、可快速接入私有文档；局限在于 retrieval score 不一定等于 generator utility，chunk 排名与 prompt 位置可能错配，generator 也可能依赖参数记忆覆盖证据。
\end{knowledgebox}

\begin{warningbox}{Frozen 不等于便宜}
每个请求仍需 embedding、search、reranking、prompt tokens 与 generation。若 top-$k$ 过大，输入 token 成本和 latency 可能超过检索收益；若 top-$k$ 太小，关键证据可能丢失。
\end{warningbox}

\teachervoice{Kiela 把 frozen RAG 称为 ``In-Context Learning only''，并反复指出它是很强的实用起点，却也是各组件独立训练的 Frankenstein system。}

\subsection{Contextualization via Retrieval：先训练哪一边}

下一张 architecture slide 用红框突出 documents、document encoder、retriever 与 query encoder，表示先让 retrieval side 适应任务。Generator 仍可冻结，训练信号则来自检索监督、最终答案 likelihood 或 reranking objective。

\lecturefigure{slide-12-contextualization-via-retrieval.jpg}{先 contextualize retrieval side 的架构路径}{Stanford 课堂视频 00:12:09}

\begin{knowledgebox}{三种常见 retriever supervision}
\begin{enumerate}
  \item \textbf{Direct relevance labels}：query--positive passage pairs 与 hard negatives；
  \item \textbf{Answer-containing heuristic}：文档是否包含答案字符串，便宜但噪声大；
  \item \textbf{Generator-derived signal}：哪份文档让正确答案 likelihood 更高，更贴近最终任务但计算更贵。
\end{enumerate}
\end{knowledgebox}

\teachervoice{课堂的 taxonomy 不是“训练或不训练”二选一，而是逐项询问 LM、query encoder、document encoder、reranker 与 index。后续所有模型都可以放回这张表定位。}

\subsection{本章小结}

RAG 架构需要同时报告 parameter updates、index refresh 与 inference behavior。Frozen RAG 用最少训练快速获得能力；更紧耦合的方法则让 retriever 或 generator 接收最终任务信号，代价是训练复杂度、索引刷新和计算成本上升。

